\section{The integer Leech shell and tail conventions}\label{app:leech}

This appendix fixes the normalization of the classical Golay
description of the Leech shell~\cite{conway1999sphere}. Let
$G\subset\F_2^{24}$ be generated by the twelve words in
Table~\ref{tab:golay}, with coordinate zero at the left. Enumeration
and binary elimination give
\begin{equation}\label{eq:golay}
 \dim G=12,\qquad G\subseteq G^\perp,\qquad
 W_G(z)=1+759z^8+2576z^{12}+759z^{16}+z^{24}.
\end{equation}
An \emph{octad} is the support of a weight-eight word of $G$.
Two distinct octads intersect in at most four places, because their
symmetric difference is a nonzero word of weight at least eight.
For $c\in G$ put $\eps_i(c)=(-1)^{c_i}$ and define
\begin{align}
 \Sigma_A&=\{4\eps e_i+4\eps'e_j:0\le i<j<24,\ \eps,\eps'\in\{\pm1\}\},
 \label{eq:shell-A}\\
 \Sigma_B&=\left\{2\sum_{i\in O}\eta_i e_i:
       O\text{ an octad},\ \eta_i\in\{\pm1\},\ \prod_{i\in O}\eta_i=1\right\},
 \label{eq:shell-B}\\
 \Sigma_C&=\left\{\sum_{i=0}^{23}\eps_i(c)e_i-4\eps_j(c)e_j:
                         c\in G,\ 0\le j<24\right\}.
 \label{eq:shell-C}
\end{align}
Set $\Sigma=\Sigma_A\cup\Sigma_B\cup\Sigma_C$.

\begin{lemma}\label{lem:shell}
This set satisfies~\eqref{eq:leech-shell}.
\end{lemma}
\begin{proof}
The absolute coordinate patterns distinguish the three families.
Support and signs recover the parameters of the first two. In the
third, the unique coordinate of absolute value three identifies $j$,
after which the signs recover $c$. Thus the cardinalities are
\[
 |\Sigma_A|=4\binom{24}{2}=1104,\quad
 |\Sigma_B|=759\cdot128=97152,\quad
 |\Sigma_C|=4096\cdot24=98304.
\]
They sum to $196560$. The respective squared norms are
$16+16$, $8\cdot4$, and $9+23$, all equal to $32$.

For two $A$-vectors, different supports give at most one contribution
of $16$, and equal supports with different signs give at most zero.
For $A$--$B$ pairs there are at most two contributions of eight;
for $A$--$C$ pairs the bound is $4(3+1)=16$.
For $B$--$B$ pairs on different octads, the intersection bound gives
at most $4\cdot4=16$. On the same octad, different even sign patterns
differ at least twice, giving at most $4(8-4)=16$.

For a $B$--$C$ pair, let $O$ be the octad and $j$ the exceptional
coordinate of the $C$-vector. If $j\notin O$, the inner product is
at most $2|O|=16$. Otherwise put $y_i=\eta_i\eps_i(c)$ on $O$.
Self-orthogonality of $G$ and the even sign condition imply
$\prod_{i\in O}y_i=1$, and the inner product is
\[
 2\sum_{i\in O}y_i-8y_j.
\]
If $y_j=1$, this is at most eight. If $y_j=-1$, at least one other
$y_i$ is negative, so the sum is at most four and the expression is
at most sixteen.

Finally consider $C$--$C$ pairs with words $c,c'$ and exceptional
coordinates $j,k$. Set $s_i=(-1)^{c_i+c'_i}$ and
$S=\sum_i s_i=24-2\wt(c+c')$. If $j=k$, the inner product is
$S+8s_j$; distinct vectors then require $c\ne c'$, giving $S\le8$.
If $j\ne k$, it is $S-4(s_j+s_k)$. This equals sixteen when
$c=c'$, and is at most $8+8$ otherwise. These six family pairs
exhaust all cases.
\end{proof}

\begin{table}[ht]
\centering\small
\caption{A binary generator for the Golay code in~\eqref{eq:golay}.
Coordinates run from zero to $23$ from left to right.}
\label{tab:golay}
\begin{tabular}{rl}
\toprule
Row & Binary word\\
\midrule
0 & \texttt{010111100101000000000001}\\
1 & \texttt{110111001010000000000010}\\
2 & \texttt{100011011101000000000100}\\
3 & \texttt{000110111110000000001000}\\
4 & \texttt{011101011100000000010000}\\
5 & \texttt{111001101001000000100000}\\
6 & \texttt{001001110111000001000000}\\
7 & \texttt{011010001111000010000000}\\
8 & \texttt{101100011011000100000000}\\
9 & \texttt{101111110000001000000000}\\
10 & \texttt{110010110011010000000000}\\
11 & \texttt{111100100110100000000000}\\
\bottomrule
\end{tabular}

\end{table}

For compatibility with the supplied tail arrays, label $V$ by
\begin{equation}\label{eq:tail-labels}
\begin{array}{c|rrr@{\qquad}c|rrr}
i&\multicolumn{3}{c}{v_i}&i&\multicolumn{3}{c}{v_i}\\\hline
0&-1&0&-1&6&-1&0&1\\
1&0&-1&-1&7&1&-1&0\\
2&-1&-1&0&8&0&-1&1\\
3&0&1&-1&9&1&1&0\\
4&-1&1&0&10&0&1&1\\
5&1&0&-1&11&1&0&1.
\end{array}
\end{equation}
The triangles in~\eqref{eq:triangles} then have labels
$\{0,7,10\}$, $\{1,6,9\}$, $\{2,3,11\}$, and $\{4,5,8\}$.
For the pure tails set
\[
 (\sigma(0),\ldots,\sigma(11))=(0,3,4,1,5,2,9,6,10,7,8,11),
 \qquad p_a=Rv_{\sigma(a)}/\sqrt2.
\]
The source rows $\widehat t_i,\widehat p_a$ in
\qpath{constructions/leech/data/tails.json}
are related to these coordinates by right multiplication by
\begin{equation}\label{eq:tail-map}
 O=\begin{pmatrix}
 \sqrt2/2&0&\sqrt2/2\\
 -\sqrt6/6&\sqrt6/3&\sqrt6/6\\
 -\sqrt3/3&-\sqrt3/3&\sqrt3/3
 \end{pmatrix},\qquad O^{\mathsf T}O=I.
\end{equation}
In row-vector notation,
$\widehat t_iO=v_i^{\mathsf T}/\sqrt6$ and
$\widehat p_aO=v_{\sigma(a)}^{\mathsf T}R^{\mathsf T}/\sqrt2$.
The exact coordinate-change test checks these identities for every
label, so the common rotation preserves both within-set and cross-set
inner products.
